\subsection{Organizzazione delle campagne}

Le campagne sono progressive e hanno obiettivi distinti. Gli Stage A--C verificano generazione dei target statici, trasferimento dello stato e comportamento della GNN statica; lo Stage D valuta il modello trajectory-aware nel solver PDE; lo Stage E1 caratterizza il solutore e la sua formulazione; lo Stage F1 confronta GDM con metodi classici di $r$-refinement. I campioni validation sono utilizzati per la selezione dei parametri classici e per le diagnostiche di sviluppo; i test held-out restano separati dalla selezione.

\paragraph{Stage A: dataset statico.}
Il dataset statico comprende 210 campioni su mesh $32\times32$: 120 training, 30 validation, 30 test interpolativi e 30 test estrapolativi. I campi sono gaussiani, con centri campionati nel dominio, ampiezza nell'intervallo $[0.8,1.2]$ e parametro $\sigma$ nell'intervallo nominale $[0.03,0.08]$; i test estrapolativi includono valori di $\sigma$ più piccoli e più grandi. I target sono accettati sulla base della riduzione del funzionale e dell'errore d'interpolazione e della validità geometrica. L'accettazione non implica stazionarietà dell'ottimizzatore.

\paragraph{Stage B--C: trasferimento e GNN statica.}
Queste campagne confrontano trasferimento puntuale $P_1$, proiezione $L^2$ e movimento ALE su un benchmark preliminare. Il target denominato nelle prime versioni ``oracle'' è ottenuto con tre iterazioni L-BFGS e viene qui chiamato \emph{baseline variazionale a basso budget}; non rappresenta un ottimo né un limite superiore. Le campagne Stage B--C sono mantenute come diagnostica dei componenti e non vengono fuse con la valutazione nativa Stage D--F1; il percorso Stage C usa il prototipo Python/SciPy descritto nell'audit della repository.

\paragraph{Stage D: benchmark PDE trajectory-aware.}
Il dataset contiene 88 campioni: 48 training, 12 validation, 12 test interpolativi e 16 test estrapolativi. Il test estrapolativo riguarda principalmente valori di $\sigma$ esterni all'intervallo di training; non rappresenta generalizzazione a PDE o geometrie differenti. Nel test F1, $\sigma$ varia da 0.02236 a 0.10339 (fattore 4.62) e $\beta$ da 0.2596 a 0.7396, mentre il vettore $\bm b$, la diffusività, il dominio e la traiettoria $\bm c(t)$ restano invariati. La variazione di $\beta$ modifica il decadimento temporale globale dell'ampiezza; non cambia il centro o la velocità prescritta della struttura. I 28 casi sono quindi un insieme parametrico della medesima famiglia MMS, non 28 problemi fisicamente indipendenti con differenti direzioni di trasporto. Il benchmark held-out usa mesh iniziale $32\times32$ (1089 nodi, 2048 triangoli), $T=0.5$, $\Delta t=0.005$, 100 passi e adattamento ogni dieci passi. Le strategie sono Fixed FEM, Stage A GNN e Stage D GNN trajectory-aware. Non è effettuato remapping esplicito; il solve dinamico usa il core nativo C++20/deal.II e la formulazione Galerkin.

\paragraph{Stage E1: verifica numerica del solver.}
La convergenza temporale è studiata su mesh $32\times32$ per $\Delta t=0.01,0.005,0.0025,0.00125$. La convergenza spaziale usa mesh $16\times16$, $32\times32$ e $64\times64$, con controllo separato della contaminazione temporale. Il Péclet locale è calcolato per elemento su mesh fissa e ALE; la SUPG è confrontata con Galerkin su un caso di riferimento. La D-GCL è verificata con test indipendenti element-wise, algebrici e free-stream. Il work--precision usa mesh fisse con $\Delta t=0.0025$, un warm-up escluso e cinque ripetizioni CPU per configurazione.

\paragraph{Stage F1: confronto held-out con $r$-refinement classico.}
La campagna confronta cinque strategie: mesh fissa, MMPDE-type Winslow, MMPDE-type arc-length, monitor Hessian-based e GDM Stage D. Tutti i metodi usano 1089 nodi, 2048 triangoli, identica connettività iniziale, $T=0.5$, $\Delta t=0.005$, adattamento ogni dieci passi e lo stesso solver nativo ALE/FEM. Non sono effettuati $h$-refinement, coarsening o remeshing. I tre parametri di intensità dei monitor classici sono selezionati soltanto sulla validation: $\alpha_W=0.1$, $\alpha_A=0.025$, $\alpha_H=0.005$. Ciascun valore è il migliore tra quattro punti discreti esplorati per il corrispondente monitor, non un ottimo continuo o una selezione esaustiva; i valori testati e gli errori validation sono riportati nell'Appendice~\ref{app:stageef}. Il test comprende 28 casi (12 interpolativi e 16 estrapolativi). Per ciascun caso e metodo è salvata una traiettoria di accuratezza; i tempi sono misurati separatamente con un warm-up e cinque ripetizioni.

\subsection{Definizione delle metriche}

L'endpoint primario è l'errore relativo $L^2$ mediato nel tempo, definito in \eqref{eq:time_average_error}. Per i confronti F1, la riduzione relativa appaiata rispetto a una baseline $b$ è
\begin{equation}
R_b(c)=100\,\frac{\overline E_{L^2}^{(b)}(c)-\overline E_{L^2}^{(GDM)}(c)}{\overline E_{L^2}^{(b)}(c)}.
\end{equation}
Valori positivi indicano un errore inferiore per GDM. L'unità statistica è il caso PDE, non la ripetizione temporale. Gli intervalli sono calcolati mediante bootstrap appaiato a livello di caso; il test inferenziale è Wilcoxon signed-rank esatto a due code e la correzione multipla usa Holm sui quattro confronti con le tre coorti (totale, interpolativa, estrapolativa). Nei confronti totali, tutti i 28 segni sono favorevoli a GDM; il corrispondente $p$ grezzo raggiunge quindi il minimo esatto $2/2^{28}=7.45\times10^{-9}$, mentre il valore Holm riportato riflette la correzione sulla famiglia di 12 test. Questo valore minimo descrive la coerenza del segno nel campione, non la grandezza dell'effetto né la variabilità tra addestramenti.

Le metriche geometriche includono minimo Jacobiano orientato, qualità minima di Escobar, aspect ratio massimo e numero di celle invertite/degenerate. Per il filtro si registrano fattore globale $\alpha$, spostamento raw e safe e slack rispetto alla soglia. La qualità geometrica viene riportata separatamente dall'accuratezza: una riduzione dell'errore non implica automaticamente un miglioramento della qualità degli elementi.

Il costo viene descritto come tempo online di implementazione. Le baseline MMPDE della F1 sono controller NumPy/Python, mentre GDM usa inferenza PyTorch; i tempi non sono pertanto un confronto di complessità algoritmica né una misura indipendente dal linguaggio e dall'implementazione. Il costo offline di training GDM non è incluso nel tempo online e va considerato separatamente in un'analisi ammortizzata.

\subsection{Discrepanza degli incrementi di massa}
La discrepanza tra gli incrementi di massa numerici e quelli della soluzione di riferimento è definita come
\begin{equation}
D_M^n=
\left|
(M_h^n-M_h^{n-1})-
(M_{\mathrm{ref}}^n-M_{\mathrm{ref}}^{n-1})
\right|,
\label{eq:mass_increment_discrepancy}
\end{equation}
dove le masse numerica e di riferimento sono
\begin{align*}
M_h^n&=\int_\Omega u_h^n\,d\bm x,\\
M_{\mathrm{ref}}^n&=\int_\Omega u^\star(\bm x,t_n)\,d\bm x.
\end{align*}
La versione relativa è normalizzata mediante il fattore
$\max(|M_{\mathrm{ref}}^0|,10^{-12})$. Questa quantità misura lo scarto tra
incrementi numerici e di riferimento; non è il residuo completo del bilancio
integrale della PDE e non viene denominata errore di conservazione. Nel caso
corrente $g_D\equiv0$: non vi è contributo convettivo di bordo associato a un
valore non omogeneo imposto, mentre il flusso diffusivo normale può essere non
nullo. Un bilancio discreto completo richiederebbe di sommare il residuo FEM
anche sui gradi di libertà di Dirichlet vincolati, interpretando tali componenti
come reazioni di bordo, e di confrontare l'identità risultante con la forzante
e i flussi. Tale residuo nodale completo non è calcolato nella campagna qui
riportata.

\subsection{Riproducibilità e limiti di validità}

La pipeline numerica Stage D--F1 usa il core C++20/deal.II nativo, mentre la gestione dei dataset, l'inferenza GNN e alcuni controller classici rimangono in Python. Le versioni dei checkpoint, della configurazione e dell'estensione nativa sono registrate negli snapshot degli esperimenti. Le conclusioni di generalizzazione restano limitate alla variazione parametrica della famiglia MMS considerata. La verifica E1 mostra che la coppia $16^2$--$32^2$ soddisfa il criterio operativo di contaminazione temporale adottato, mentre il livello $64^2$ non è spazialmente isolato. La campagna F1 è eseguita a $\Delta t=0.005$: non viene pertanto attribuita al solo errore spaziale la differenza di accuratezza tra le strategie, né la verifica temporale E1 su mesh fissa è utilizzata come stima dell'errore temporale specifico del benchmark mobile.
